% BHU PYQ -- Manually transcribed & error-corrected
% Paper: GLB-401: Palaeontology and Stratigraphy
% Exam: B.Sc. (Hons.) Semester IV Examination, 2013-14

\documentclass[11pt,a4paper]{article}
\usepackage[a4paper,top=2.5cm,bottom=2.5cm,left=2.8cm,right=2.8cm,headheight=14pt]{geometry}
\usepackage{amsmath,amssymb}
\usepackage[T1]{fontenc}
\usepackage[utf8]{inputenc}
\usepackage{lmodern,microtype,fancyhdr,enumitem,hyperref,lastpage,parskip}

\hypersetup{colorlinks=true,urlcolor=black,linkcolor=black,
  pdftitle={GLB-401 Palaeontology and Stratigraphy -- BHU B.Sc. Sem IV 2013-14}}
\pagestyle{fancy}\fancyhf{}
\renewcommand{\headrulewidth}{0.4pt}\renewcommand{\footrulewidth}{0.4pt}
\fancyhead[L]{\small   \textit{Banaras Hindu University}}
\fancyhead[R]{\small   \textit{GLB-401: Palaeontology \& Stratigraphy}}
\fancyfoot[C]{\small Page  \thepage\ of \pageref{LastPage}}


\newlist{parts}{enumerate}{2}
\setlist[parts,1]{label=(\alph*),leftmargin=*,itemsep=5pt,topsep=3pt}

\newlist{romanparts}{enumerate}{2}
\setlist[romanparts,1]{label=(\roman*),leftmargin=*,itemsep=5pt,topsep=3pt}

\newcommand{\pts}[1]{\hfill   \textit{[#1]}}


\begin{document}

\begin{center}
  {\large\bfseries BANARAS HINDU UNIVERSITY}\\[2pt]
  {
\normalsize B.Sc. (Hons.) --- Semester IV Examination, 2013-14}\\[6pt]
  \rule{0.8\linewidth}{0.5pt}\\[6pt]
  {\large\bfseries Paper: GLB-401 --- Palaeontology and Stratigraphy}\\[4pt]
  {
\normalsize Time: 3 Hours \hfill Maximum Marks: 70}\\[2pt]
  {\small Ref. No.: BS/Sem IV/}
\end{center}
\rule{\linewidth}{0.4pt}
\smallskip



\noindent   \textit{Instructions: Answer   \textbf{five} questions in all, selecting at least   \textbf{two} each from Sections A and B, and Section-C which is compulsory. All questions carry equal marks.}
\bigskip

\bigskip


\noindent {\large\bfseries SECTION --- A}
\rule{\linewidth}{0.4pt}\smallskip




\noindent   \textbf{Question 1.} Define Palaeontology and describe its relationship with other branches of geology. \pts{14}

\medskip


\noindent   \textbf{Question 2.} Briefly discuss   \textbf{any two} of the following: \pts{$2  \times 7 = 14$}

\begin{parts}
  \item Essential conditions of fossilization
  \item Origin of life
  \item Various terms used to describe the bathymetric distribution of organisms
\end{parts}

\medskip


\noindent   \textbf{Question 3.} Describe the morphological features of Trilobites. Support your answer with neat, neatly labeled diagrams. \pts{14}

\medskip


\noindent   \textbf{Question 4.} Write short notes on   \textbf{any two} of the following: \pts{$2  \times 7 = 14$}

\begin{parts}
  \item Incompleteness of the fossil record
  \item Geological Time Scale
  \item Economic importance of the Cuddapah Group
\end{parts}

\bigskip


\noindent {\large\bfseries SECTION --- B}
\rule{\linewidth}{0.4pt}\smallskip




\noindent   \textbf{Question 5.} Describe in detail the various principles of stratigraphy. Illustrate your answer with neat diagrams. \pts{14}

\medskip


\noindent   \textbf{Question 6.} Describe in detail the distribution, classification, and economic importance of the Dharwar Supergroup. Briefly discuss what kind of geological information is derived from these rocks. \pts{14}

\medskip


\noindent   \textbf{Question 7.} Write short notes on   \textbf{any two} of the following: \pts{$2  \times 7 = 14$}

\begin{parts}
  \item In what way stratigraphy requires inputs from all branches of geology
  \item What is the significance of the Eparchaean Unconformity
  \item Why unconformities are often marked by conglomerates
\end{parts}

\bigskip


\noindent {\large\bfseries SECTION --- C}
\rule{\linewidth}{0.4pt}\smallskip




\noindent   \textbf{Question 8.} Distinguish between or answer in 1 to 3 sentences: \pts{$7  \times 2 = 14$}

\begin{romanparts}
  \item Fossils and pseudofossils
  \item Species and subspecies
  \item Sargur Group
  \item \emph{Chondrites} and dendrites
  \item Diamondiferous conglomerate
  \item Lithostratigraphic units
  \item Apical disc in Echinoids
\end{romanparts}

\vfill\rule{\linewidth}{0.4pt}

\begin{center}\small
\href{https://bhu-exam.vercel.app/}{Click here to visit our website}\\[4pt]
\href{https://me-aryan.vercel.app/}{Click here to visit our contributor}\\[4pt]
\href{https://quashan-me.vercel.app/}{Click here to visit our contributor}
\end{center}
\end{document}
